\chapter{Coronary Angiogram}

\section*{Definition and Overview}
A coronary angiogram (coronary angiography) is an X-ray test that shows the inside of the coronary arteries, the blood vessels that supply the heart muscle. A thin, flexible tube called a catheter is inserted into an artery, usually in the wrist or groin, and guided to the heart under X-ray guidance. Contrast dye is injected through the catheter, and X-ray images show whether the coronary arteries are narrowed or blocked by fatty deposits (plaque). The test is the standard method for diagnosing coronary artery disease and is used to guide treatment such as stenting or bypass surgery.

\section*{When to Seek Medical Care}
A coronary angiogram is performed in a hospital cardiac catheterisation laboratory, usually as a planned (elective) procedure or urgently after a heart attack or unstable angina. Before the test, your doctor will discuss the reasons for it and the risks.

\textbf{Call 000 (or local emergency services) for:}
\begin{itemize}
    \item Chest pain, pressure, or discomfort that is new, worsening, or not relieved by rest
    \item Pain spreading to the arm, jaw, neck, or back
    \item Shortness of breath, sweating, nausea, or feeling faint with chest pain
    \item A heart attack emergency, or symptoms that occur after a recent heart procedure
\end{itemize}

\section*{Epidemiology}
Coronary artery disease is one of the leading causes of death and disability worldwide. In Australia, it remains a major cause of illness and death, and coronary angiography is among the most commonly performed hospital procedures. The test is used to evaluate people with angina, heart attacks, abnormal stress tests, or heart failure of uncertain cause. Rates of angiography have fallen slightly in recent years in some countries as non-invasive tests have improved, but it remains the reference standard for defining the anatomy of the coronary arteries.

\section*{Aetiology and Pathophysiology}
The coronary arteries narrow because of atherosclerosis, the gradual buildup of fatty plaque within the artery walls. A coronary angiogram is performed to answer specific questions about these vessels.

\begin{itemize}
    \item \textbf{Atherosclerosis:} cholesterol, fat, and inflammatory cells accumulate in the artery wall, narrowing the lumen
    \item \textbf{Plaque rupture:} a plaque can tear, causing a blood clot that suddenly blocks the artery and causes a heart attack
    \item \textbf{Stable angina:} fixed narrowing that limits blood flow during exertion
    \item \textbf{Acute coronary syndromes:} unstable angina and heart attack from partial or complete blockage
    \item \textbf{Assessment of prior treatment:} checking grafts or stents after bypass surgery or angioplasty
    \item \textbf{Pre-operative assessment:} evaluating coronary anatomy before certain heart valve or other surgeries
\end{itemize}

\section*{Clinical Features}
Angiography is not a treatment; it is a diagnostic test. The symptoms that lead to it are those of possible coronary disease.

\begin{itemize}
    \item Chest pain or discomfort, often described as pressure, squeezing, or heaviness
    \item Pain or discomfort spreading to the arm, jaw, neck, or back
    \item Shortness of breath, especially with exertion
    \item Fatigue, nausea, or sweating with exertion
    \item Palpitations or an irregular heartbeat
    \item Abnormal results on an exercise stress test, ECG, or CT coronary calcium score
\end{itemize}

\section*{Differential Diagnosis}
Before angiography, other causes of chest pain are usually considered and excluded.

\begin{itemize}
    \item Gastro-oesophageal reflux or oesophagitis
    \item Musculoskeletal chest pain, including costochondritis
    \item Anxiety and panic attacks
    \item Pulmonary embolism, pneumonia, or pleurisy
    \item Aortic dissection, a tearing of the wall of the aorta
    \item Pericarditis, inflammation of the sac around the heart
\end{itemize}

\section*{Diagnosis and Investigations}
The procedure itself is carefully planned and monitored.

\begin{itemize}
    \item \textbf{Pre-procedure assessment:} blood tests, ECG, and review of kidney function and allergies
    \item \textbf{Catheter insertion:} usually via the radial artery in the wrist, or the femoral artery in the groin, under local anaesthetic
    \item \textbf{Contrast injection:} dye is injected and X-ray images (cineangiograms) record blood flow through each coronary artery
    \item \textbf{Pressure measurements:} pressures inside the heart chambers may be recorded
    \item \textbf{Fractional flow reserve (FFR):} a wire measures whether a narrowing significantly limits blood flow
    \item \textbf{Intravascular imaging:} ultrasound or optical coherence tomography can characterise plaques
    \item \textbf{Immediate treatment:} if a significant blockage is found, angioplasty and stenting can be performed in the same session
\end{itemize}

\section*{Management}
The results of the angiogram guide further treatment.

\begin{itemize}
    \item \textbf{No significant disease:} risk factor control, lifestyle measures, and regular review
    \item \textbf{Medical therapy:} antiplatelet drugs, statins, and medicines for angina or heart failure
    \item \textbf{Percutaneous coronary intervention (PCI):} balloon angioplasty and stenting of significant narrowings
    \item \textbf{Coronary artery bypass grafting (CABG):} surgery for extensive or complex disease
    \item \textbf{Revascularisation planning:} the heart team discusses the best option based on the anatomy
    \item \textbf{Risk factor management:} smoking cessation, blood pressure control, cholesterol lowering, and diabetes management
    \item \textbf{Cardiac rehabilitation:} exercise, education, and support after the procedure
\end{itemize}

\section*{Complications}
A coronary angiogram is generally safe, but it is an invasive test and carries small risks.

\begin{itemize}
    \item Bleeding, bruising, or a haematoma at the catheter site
    \item Damage to the artery at the puncture site, occasionally needing repair
    \item An allergic reaction to the contrast dye
    \item Kidney injury from the contrast, especially in people with pre-existing kidney disease
    \item Irregular heart rhythms during the procedure
    \item Stroke, heart attack, or damage to a coronary artery (very rare)
    \item Infection at the puncture site
    \item Radiation exposure from X-rays
\end{itemize}

\section*{Prognosis and Outcomes}
For most people, a coronary angiogram is a safe procedure with a rapid recovery, and most can go home the same day or the next morning. The main value of the test is that it clarifies the diagnosis and directs treatment, and people with significant disease can be treated immediately with stenting or referred for bypass surgery, which improves symptoms and, in many cases, survival. People with normal or mildly abnormal results are usually reassured and managed with risk factor control and lifestyle measures.

\section*{Prevention and Self-Care}
\begin{itemize}
    \item Follow the preparation instructions, including fasting and medication adjustments
    \item Tell the team about allergies, kidney problems, diabetes, and all medicines
    \item After the procedure, keep the puncture site clean and avoid heavy lifting for a few days
    \item Drink fluids as advised to help clear the contrast dye
    \item Attend follow-up and take prescribed medicines exactly as directed
    \item Manage risk factors: do not smoke, control blood pressure, cholesterol, and blood sugar, and stay active
    \item Attend cardiac rehabilitation if it is offered
\end{itemize}

\section*{Sources and Further Reading}
The following reputable sources provide further information for healthcare students and patients seeking reliable, current advice about coronary angiography.

\begin{itemize}
    \item NHS. \textit{Coronary angiogram: what happens and risks.} https://www.nhs.uk/conditions/coronary-angiogram/
    \item Mayo Clinic. \textit{Coronary angiogram: overview and what to expect.} https://www.mayoclinic.org/tests-procedures/coronary-angiogram/
    \item Heart Foundation Australia. \textit{Coronary angiogram.} https://www.heartfoundation.org.au/
    \item MedlinePlus. \textit{Coronary angiography health information.} https://medlineplus.gov/ency/article/003876.htm
    \item MSD Manual. \textit{Coronary angiography: professional reference.} https://www.msdmanuals.com/
    \item American Heart Association. \textit{Cardiac catheterisation and angiography.} https://www.heart.org/
\end{itemize}
